% ============================================================================
% บทที่ 9 — ปัญหาที่พบบ่อยและรายการตรวจงาน
% ============================================================================
\mychapter{ปัญหาที่พบบ่อยและรายการตรวจงาน}

\section{ตารางปัญหาที่พบบ่อย}

\begin{longtable}{L{4.5cm}L{6cm}L{4cm}}
\toprule
\textbf{อาการ} & \textbf{การแก้ไข} & \textbf{ขั้นตอนที่เกี่ยวข้อง} \\
\midrule
ปุ่มนำเข้าไม่ทำงาน & ตรวจสอบว่ามีไฟล์ \texttt{จัดอัตโนมัติ} อย่างน้อยหนึ่งไฟล์ & บทที่ 3 \\
ช่วงเวลาสอบไม่ถูกต้อง & ตรวจสอบวันที่เริ่ม/สิ้นสุดของกลางภาคและปลายภาค & บทที่ 3 \\
รหัสวิชาไม่พบในเงื่อนไข & ตรวจสอบเลขศูนย์นำหน้าของรหัสวิชา & บทที่ 2 \\
วิชาแก้ไขเวลาไม่ได้ & วิชามีเวลาจากไฟล์ต้นทางหรือถูกบล็อก & บทที่ 5 \\
การแก้ไขเวลาถูกปฏิเสธ & เวลาทับซ้อนกับวิชาอื่นในกลุ่มเดียวกัน & บทที่ 5 \\
วิชาไม่ถูกจัดตาราง & ปัญหาข้อมูลต้นทาง เช่น กลุ่มนักศึกษาขาดหาย & บทที่ 6 \\
จัดได้ไม่ครบแม้ไม่มีข้อผิดพลาด & เพิ่มระดับความพยายามแล้วจัดใหม่ ระบบหาคำตอบได้ดีขึ้นแต่ไม่ได้พิสูจน์ว่าตารางเป็นไปไม่ได้ & บทที่ 3 \\
การจัดตารางใหม่ลบการแก้ไข & บันทึกโครงการขึ้นคลาวด์ก่อนจัดตารางใหม่ & บทที่ 5 \\
การบันทึกขึ้นคลาวด์ไม่ทำงาน & ตรวจสอบว่ามีบริการคลาวด์ที่ตั้งค่าไว้ & บทที่ 8 \\
เกิดข้อขัดแย้งฉบับ & ใช้ตัวเลือกดาวน์โหลด/แก้ต่อ/โหลดจากคลาวด์ & บทที่ 8 \\
\bottomrule
\end{longtable}

\section{คู่มือและตัวอย่างในระบบ}

หน้า \textbf{วิธีใช้งาน คู่มือและตัวอย่าง} ประกอบด้วยแท็บ
\texttt{วิธีใช้งาน} (คำแนะนำการใช้งานทั่วไป) \texttt{รูปแบบข้อมูล}
(โครงสร้างไฟล์ข้อมูล) \texttt{ตัวอย่างข้อมูล} และ \texttt{แก้ปัญหา}
(วิธีแก้ปัญหาที่พบบ่อย) สามารถคัดลอกตัวอย่างข้อมูลจากแท็บ
\texttt{ตัวอย่างข้อมูล} ไปใช้เป็นจุดเริ่มต้นได้

\section{รายการตรวจสอบก่อนการแชร์}

ก่อนส่งลิงก์แชร์ให้ผู้อื่น ควรตรวจสอบรายการต่อไปนี้:

\begin{itemize}
  \item บันทึกฉบับล่าสุดของโครงการแล้วหรือยัง
  \item เลือกสิทธิ์ที่เหมาะสมสำหรับผู้รับแล้วหรือยัง
  \item เลือกฉบับ (\texttt{ตามฉบับล่าสุด} หรือ \texttt{ตรึงฉบับที่บันทึกไว้})
        ที่ต้องการแล้วหรือยัง
  \item กำหนดอายุของลิงก์ตามความเหมาะสมแล้วหรือยัง
  \item ตรวจสอบว่าการแก้ไขที่ต้องการให้ผู้รับเห็นถูกบันทึกแล้วหรือยัง
\end{itemize}

\section{รายการตรวจสอบหลังการเปิดงานต่อ}

หลังเปิดงานต่อ ควรตรวจสอบรายการต่อไปนี้:

\begin{itemize}
  \item ข้อมูลโครงการที่เปิดถูกต้องหรือไม่
  \item การแก้ไขเวลาที่ยังต้องการอยู่ครบถ้วนหรือไม่
  \item จำนวนวิชาที่ยังไม่ได้จัดเป็นเท่าไร
  \item ผลการตรวจสอบมีข้อผิดพลาดหรือคำเตือนหรือไม่
\end{itemize}

\seealso{บทที่ 2 สำหรับการเตรียมข้อมูล และบทที่ 6 สำหรับการตรวจสอบปัญหา}
